% Section 9.1, from lectures/L09/appendix/a_sleeping.md.

\section{Sleeping and Waiting}
\label{c9:app:a}

The other half of Chapter~\ref{c8:ch}. That chapter was about the half the hardware drives: a line
goes high and a handler runs. This one is about the half the software drives: a program has asked
for data that does not exist yet, and something has to happen to it in the meantime.

\Secref{c9:app:b} is time, delays and timers.

The idea to carry: \textbf{sleeping is not waiting.} A sleeping task is off the run queue entirely
and costs nothing; a spinning task is on it and costs a CPU. The whole of this section is machinery
for getting into the first state safely.

\subsection{What sleeping is}
\label{c9:app:a1}

A task has a state. The ones that matter:

\begin{booktable}{@{}>{\hsize=0.80\hsize}Lll>{\hsize=1.20\hsize}L@{}}
  \thead{State} & \thead{\code{ps} shows} & \thead{On the run queue} & \thead{Woken by} \\
  \midrule
  \code{TASK_RUNNING} & \code{R} & Yes & It is already running or ready \\
  \code{TASK_INTERRUPTIBLE} & \code{S} & \textbf{No} & An explicit wake, or a signal \\
  \code{TASK_UNINTERRUPTIBLE} & \code{D} & \textbf{No} & An explicit wake only \\
\end{booktable}

\noindent To sleep, a task sets its state to something other than \code{TASK_RUNNING} and calls
\code{schedule()}. The scheduler then does not consider it again until something sets it back to
\code{TASK_RUNNING}. That is all a wait queue is built out of.

\textbf{Use \code{TASK_INTERRUPTIBLE} in a driver, essentially always.} The uninterruptible state
exists for short, uninterruptible hardware operations, and a task stuck in \code{D} cannot be
killed, not even with signal 9, which is \secref{c2:app:a5}'s observation arriving with a cause.

\subsection{The wrong version, and why it is wrong every time}
\label{c9:app:a2}

The obvious implementation of ``wait until there is data'':

\begin{ccode}
/* WRONG */
while (fifo_empty())
        schedule();
\end{ccode}

\noindent Two things are wrong. It never sets its state, so it stays \code{TASK_RUNNING} and the
scheduler keeps picking it: this is a busy-wait with extra steps, and it burns a CPU doing nothing.
And it is on no wait queue, so even if it did sleep, no wake-up could ever find it.

The obvious fix is worse, because it looks right. Take it that the reader has already joined the
queue the handler wakes; the snippet leaves that line out, because what is wrong is elsewhere:

\begin{ccode}
/* WRONG, and subtly */
if (fifo_empty()) {
        set_current_state(TASK_INTERRUPTIBLE);
        schedule();
}
\end{ccode}

\noindent Consider the interleaving. The reader evaluates \code{fifo_empty()} and finds it true.
\textbf{The interrupt fires here.} The handler adds data and calls \code{wake_up_interruptible},
which finds nothing asleep on the queue and does nothing at all. The reader then sets its state and
calls \code{schedule()}, and sleeps waiting for a wakeup that has already happened: until the next
one, if there is a next one, and forever if there is not.

\textbf{This is the lost wakeup, and it is not a rare race.} The window is two instructions wide,
but with an interrupt arriving every millisecond it is hit within seconds. It is not the kind of bug
that waits for production; it is the kind that fails in your first test and then works on the
second, which is worse.

\lecturefigure{lectures/L09/appendix/images/lost_wakeup.png}
  {A sequence diagram with a lane for the reader and a lane for the handler. The reader tests
   \code{fifo_empty} and finds it true; the handler then pushes a sample and calls \code{wake_up},
   which finds nothing asleep to wake; the reader then sets its state and calls \code{schedule}, and
   sleeps until the next wake, if one ever comes.}
  {fig:lost_wakeup}

\subsection{What the macro actually does}
\label{c9:app:a3}

\code{wait_event_interruptible} is the correct version, and reading its expansion is worth more than
memorising its usage. From \file{include/linux/wait.h}, abridged:

\begin{ccode}
for (;;) {
        long __int = prepare_to_wait_event(&wq_head, &__wq_entry, state);

        if (condition)
                break;

        if (___wait_is_interruptible(state) && __int) {
                __ret = __int;
                goto __out;
        }

        schedule();
}
finish_wait(&wq_head, &__wq_entry);
\end{ccode}

\noindent Three things close the window in \secref{c9:app:a2}, in this order:
\begin{enumerate}
  \item \textbf{\code{prepare_to_wait_event} first.} It adds the task to the queue \emph{and} sets
    its state, before the condition is looked at. From that moment a \code{wake_up} cannot be
    missed: it will find this task on the queue and set it runnable.
  \item \textbf{The condition is tested after being queued.} If the data arrived during step 1, the
    test sees it and breaks out without ever sleeping.
  \item \textbf{It is a loop.} After \code{schedule()} returns, the condition is tested again.
\end{enumerate}
Step 3 is why the macro takes a \emph{condition} and not a flag. \textbf{A wake does not mean the
condition is true.} Several readers may be woken and the first to run may take all the data; a wake
may be spurious. The sleeper must re-check, and the loop does it for you.

Two consequences for how you write the condition: it is evaluated \textbf{many times}, so it must be
cheap and free of side effects, and it is evaluated \textbf{from the waiting task's context}, so it
must be safe to call there.

\subsection{Using it}
\label{c9:app:a4}

\begin{ccode}
static DECLARE_WAIT_QUEUE_HEAD(qa_readq);

/* the reader */
for (;;) {
        n = take_the_data();        /* under the lock; 0 if there was none */
        if (n)
                break;
        if (file->f_flags & O_NONBLOCK)
                return -EAGAIN;
        if (wait_event_interruptible(qa_readq, !fifo_empty()))
                return -ERESTARTSYS;
}

/* the interrupt handler */
if (got_data)
        wake_up_interruptible(&qa_readq);
\end{ccode}

\noindent\textbf{The reader loops for the same reason the macro does.} A wake says there was data
when the waker looked; with two readers, the other may have taken it by the time this one runs. A
reader that takes nothing goes back to waiting, because returning 0 would tell its caller the
device had finished.

\begin{narrowtable}{@{}ll@{}}
  \thead{Function} & \thead{Wakes} \\
  \midrule
  \code{wake_up_interruptible} & Tasks in \code{TASK_INTERRUPTIBLE} \\
  \code{wake_up} & Those and \code{TASK_UNINTERRUPTIBLE} \\
  \code{wake_up_interruptible_all} & The same, and every exclusive waiter, not one \\
\end{narrowtable}

\noindent An \emph{exclusive} waiter is one that asked, with
\code{wait_event_interruptible_exclusive}, to be woken on its own. The first two wake every other
waiter on the queue and at most one exclusive one.

\textbf{Waking is safe from interrupt context} and does not sleep, which is what makes
this the natural join between Chapter~\ref{c8:ch}'s handler and this chapter's reader.

\textbf{A name to avoid.} Calling the queue \code{readq} does not compile: the kernel already has
\code{readq}, the 64-bit MMIO accessor from \secref{c6:app:b4}, and the error is \code{'readq'
redeclared as different kind of symbol}. The kernel namespace is flat and short names are taken;
prefix yours with the driver's name.

\subsection{Blocking \code{read()}, properly}
\label{c9:app:a5}

Four cases, and a driver has to get all four right:

\begin{narrowtable}{@{}ll@{}}
  \thead{Situation} & \thead{Return} \\
  \midrule
  Data available & The bytes, and how many \\
  Empty, blocking & Sleep until there is data \\
  Empty, \code{O_NONBLOCK} & \textbf{\code{-EAGAIN}} \\
  Signal arrives while asleep & \textbf{\code{-ERESTARTSYS}} \\
\end{narrowtable}

\noindent \textbf{Never 0.} Zero means end of file, which is Chapter~\ref{c5:ch}'s subject and the
reason that chapter returned \code{-EAGAIN} unconditionally. This chapter upgrades that:
\code{-EAGAIN} is now the answer only for a caller who asked for it with \code{O_NONBLOCK}.

\subsubsection{\code{-ERESTARTSYS}}
\label{c9:app:a:-erestartsys}

\code{wait_event_interruptible} returns non-zero when a signal arrived. The driver returns
\code{-ERESTARTSYS}, which never reaches userspace: the kernel sees it and either restarts the
system call transparently or converts it to \code{-EINTR}, depending on how the signal handler was
installed.

\textbf{Getting this wrong makes your device un-interruptible with Ctrl-C}, and users notice. A
driver that returns \code{-EINTR} directly is close but takes the restart decision away from the
kernel; one that ignores the return value and carries on has a reader that cannot be killed.

\subsubsection{Toggling \code{O_NONBLOCK}}
\label{c9:app:a:toggling-o_nonblock}

Worth knowing because it catches people writing tests. A descriptor's blocking behaviour is a
property of the open file, changeable at run time:

\begin{ccode}
int flags = fcntl(fd, F_GETFL, 0);
fcntl(fd, F_SETFL, flags | O_NONBLOCK);
\end{ccode}

\noindent The course's own tester needs this: draining a blocking descriptor attached to a device
that produces a sample every millisecond never terminates, because there is always more data a
millisecond from now.

\subsection{\code{poll}, and what it buys}
\label{c9:app:a6}

Two lines of driver code, and every multiplexing interface in userspace works on your device.

\begin{ccode}
static __poll_t qa_poll(struct file *file, struct poll_table_struct *wait)
{
        poll_wait(file, &qa_readq, wait);
        return fifo_empty() ? 0 : (EPOLLIN | EPOLLRDNORM);
}
\end{ccode}

\noindent \textbf{\code{poll_wait} does not wait.} It registers the caller's interest in your wait
queue and returns immediately; the sleeping is done by the \code{select}/\code{poll}/\code{epoll}
machinery above you, on all the descriptors at once. Your function is then called a second time
after any wake, to ask again.

So the function must do exactly two things: register on every queue that could make it ready, and
return a mask describing readiness \emph{now}.

\begin{booktable}{@{}lL@{}}
  \thead{Bit} & \thead{Means} \\
  \midrule
  \code{EPOLLIN} & Readable \\
  \code{EPOLLRDNORM} & Readable, normal data. Conventionally set with \code{EPOLLIN} \\
  \code{EPOLLOUT} & Writable \\
  \code{EPOLLERR} & An error \\
  \code{EPOLLHUP} & Hung up \\
\end{booktable}

\noindent What it buys is that \textbf{one thread can wait on your device and a socket and a timer
together}, which is the whole architecture of most userspace daemons. Without \code{poll}, a program
that wants to watch your device and anything else needs a thread per descriptor.

Measured on the target: with the device running, \code{select} reports the descriptor readable
within a second; with the device stopped, \code{select} given a 300 ms timeout returns 0 after
\textbf{305 ms}.

\subsection{Completions}
\label{c9:app:a7}

For the common special case of ``wait for this one thing to finish once''.

\begin{ccode}
static DECLARE_COMPLETION(setup_done);

wait_for_completion(&setup_done);       /* the waiter */
complete(&setup_done);                  /* whoever finished */
\end{ccode}

\noindent A completion is a wait queue plus a ``has it happened'' flag, so it has no lost-wakeup
problem to solve: completing before anyone waits is fine, and the waiter returns immediately.

Use a completion for a one-shot event (initialisation finished, a DMA transfer completed) and a wait
queue for a recurring condition (data is available). Choosing a wait queue for a one-shot means
writing the flag yourself, which is where the lost wakeup comes back.

\code{wait_for_completion_interruptible} and \code{wait_for_completion_timeout} exist and are
usually what you want in a driver, for the same reasons as their \code{wait_event} equivalents.

\subsection{What may sleep, and where}
\label{c9:app:a8}

The rules from \secref{c7:app:a7} and \secref{c8:app:a7}, stated once more because this is the chapter
where you deliberately sleep:

\begin{narrowtable}{@{}ll@{}}
  \thead{Context} & \thead{May sleep} \\
  \midrule
  A system call, so \code{read} & \textbf{Yes} \\
  A threaded IRQ handler & \textbf{Yes} \\
  A workqueue & \textbf{Yes} \\
  A hard IRQ handler & No \\
  Holding a spinlock & No \\
  A timer callback & No \\
\end{narrowtable}

\noindent The one that catches people in this chapter: \textbf{you may not sleep while holding a
spinlock}, and \code{wait_event_interruptible} sleeps. So the condition is checked under the lock,
the lock is dropped, and the wait happens outside it. Doing it the other way round deadlocks the
machine, and \code{CONFIG_DEBUG_ATOMIC_SLEEP} turns that into a \code{BUG: sleeping function called
from invalid context} rather than a hang.
